\documentclass[11pt]{article}
\usepackage[margin=1in]{geometry}
\usepackage{amsmath,amssymb,amsthm}
\usepackage{graphicx}
\usepackage{booktabs}
\usepackage{hyperref}
\usepackage{xcolor}
\usepackage{natbib}
\usepackage{subcaption}
\usepackage{multirow}
\usepackage{longtable}

\newcommand{\citeneeded}[1]{\textcolor{red}{[CITE: #1]}}
\newcommand{\todo}[1]{\textcolor{orange}{[TODO: #1]}}

\title{Most knockout epistasis is generated by dynamics, not molecular wiring}

\author{
  Elliot Tower \\
  \texttt{elliot@elliottower.ai}
}

\begin{document}
\maketitle

\begin{abstract}
Nonlinear dynamics generate phenotypic epistasis from molecular
architectures that may contain no epistatic interactions themselves.
Prior work established this gap qualitatively; we quantify it
order-by-order across 27 published Boolean gene regulatory networks
(7--17 genes). For each network, we compute the complete
Walsh--Hadamard interaction spectrum at two levels: the local update
rules and the attractor-level knockout landscape from exhaustive
$2^n$ coalition sweeps. Comparing the two spectra reveals a signed
\emph{composition gap} spanning three orders of magnitude ($< 1$ to
$+83$ percentage points of order-3+ energy). Creation of higher-order
interactions is the dominant mode: 18 of 24 non-null networks show
positive gaps, with a median of $+8.5$~pp. Pairwise Spearman
correlation between local and global interaction magnitudes is positive
on average (median $\rho = 0.27$, significant in 16/27 networks) but
negative in five, demonstrating that molecular wiring is a noisy and
sometimes misleading predictor of knockout epistasis. Limit-cycle prevalence correlates negatively with the composition gap
(discovered on 6 networks, same direction on 21 held-out networks:
$\rho = -0.37$, $p = 0.10$), though counterexamples in both
directions preclude a clean threshold. Pre-registered structural
predictions on 21 held-out networks---including a feedback loop ratio
hypothesis---perform at chance on gap direction (10/21, 48\%),
establishing that the composition gap is not reducible to simple
topological features of the regulatory graph. The complete Walsh spectra, coalition value functions, and
pre-registration artifacts are released for reuse.
\end{abstract}


%% ====================================================================
\section{Introduction}
\label{sec:intro}
%% ====================================================================

Epistasis---the dependence of phenotype on gene combinations rather
than individual effects---is among the oldest measurement problems in
quantitative genetics \citeneeded{Fisher 1918, Bateson 1909}. A
persistent difficulty is the gap between molecular and statistical
epistasis: the interactions encoded in regulatory logic need not
predict the interactions observed in knockout experiments
\citeneeded{Phillips 2008 Nat Rev Genet}. Gjuvsland et al.\ showed
that nonlinear dynamics generically produce statistical epistasis from
gene regulatory architectures, regardless of whether the underlying
molecular rules contain epistatic terms \citeneeded{Gjuvsland 2007
Genetics}. Subsequent work confirmed this in synthetic three-gene
circuits, where pairwise and triplet epistasis depended on environment
and switched sign across conditions \citeneeded{Baier 2023 Sci Adv}.

These results establish that a gap exists. They do not quantify it.
The gap has no known magnitude, no known sign structure, and no
established dependence on network architecture. One reason is
computational: exhaustive knockout measurement requires $2^n$ subset
evaluations, which is infeasible at genomic scale. In Boolean gene
regulatory networks of modest size ($n \leq 17$), however, exhaustive
evaluation is tractable.

We exploit this tractability across 27 published Boolean networks
spanning cell-cycle control, apoptosis, differentiation, DNA repair,
and developmental patterning. For each network, we compute two
complete interaction spectra using the Walsh--Hadamard transform: one
from the local update rules and one from the attractor-level value
function. The two spectra are compared order-by-order, yielding a
signed composition gap that measures how much higher-order structure
dynamics create or destroy.

Three findings emerge. First, creation is the dominant mode: dynamics
more often generate higher-order interactions than destroy them (18/24
non-null networks). Second, the molecular wiring diagram is a noisy
and sometimes misleading predictor of knockout epistasis---positive on
average but negative in five networks. Third, pre-registered
structural predictions perform at chance on gap direction (48\%),
establishing that the composition gap is not reducible to network
topology.


%% ====================================================================
\section{Methods}
\label{sec:methods}
%% ====================================================================

\subsection{Boolean network models}
\label{sec:models}

Twenty-seven published Boolean gene regulatory networks were collected
from the Cell Collective repository \citeneeded{Helikar 2012}, the
PyBoolNet model repository \citeneeded{Klarner 2017}, the Biodivine
Boolean Models repository, and primary literature. Selection criteria:
(i)~$n \leq 18$ nodes to permit exhaustive $2^n$ coalition sweeps
within a day of CPU time, (ii)~synchronous update semantics used or
compatible with the original publication, and (iii)~a biologically
interpretable output node. The networks span cell-cycle control
(mammalian, yeast, \emph{Drosophila}, \emph{Arabidopsis}), apoptosis,
hematopoietic differentiation, DNA repair, developmental patterning,
gene regulation, and signal transduction (Table~\ref{tab:all_results}).

To guard against overfitting structural hypotheses to results, the
models were analyzed in two batches. Batch~1 (6~networks) was used to
develop the analysis pipeline and formulate structural predictions.
Batch~2 (21~networks) was analyzed blind: structural predictions were
SHA-frozen before any coalition sweep ran
(Section~\ref{sec:blind}).


\subsection{Local rule Fourier extraction}
\label{sec:local-fourier}

Each gene $k$ has a Boolean update rule $f_k\colon \{0,1\}^{r_k}
\to \{0,1\}$, where $r_k$ is the number of regulators. We compute
the Walsh--Hadamard transform of the truth table:
%
\begin{equation}
\label{eq:local-wht}
w_T^{(k)} = \frac{1}{2^{r_k}} \sum_{x \in \{0,1\}^{r_k}} f_k(x)\,(-1)^{|x \cap T|},
\quad T \subseteq \{1, \ldots, r_k\}.
\end{equation}
%
A nonzero coefficient $w_T^{(k)}$ at $|T| = m$ indicates a genuine
$m$-way interaction among the regulators of gene $k$,
representation-independent. For each unordered pair $\{i,j\}$ of
network genes, the \emph{local pairwise magnitude} is:
%
\begin{equation}
\label{eq:local-pair}
\ell_{ij} = \sum_{k:\, i,j \in \mathrm{reg}(k)} |w_{\{i,j\}}^{(k)}|,
\end{equation}
%
summing over all genes $k$ whose update rule depends on both $i$ and
$j$. The \emph{local energy spectrum} reports the fraction of total
squared-coefficient energy at each interaction order~$m$:
%
\begin{equation}
\label{eq:local-spectrum}
E_m^{\mathrm{local}} = \frac{\sum_k \sum_{|T|=m} \bigl(w_T^{(k)}\bigr)^2}
{\sum_k \sum_T \bigl(w_T^{(k)}\bigr)^2}.
\end{equation}


\subsection{Coalition sweep and attractor-level value function}
\label{sec:coalition-sweep}

For each of $2^n$ coalitions $S \subseteq \{1,\ldots,n\}$, genes not
in $S$ are clamped to~0 (knockout). The remaining genes evolve under
synchronous Boolean dynamics from $N = 512$ uniformly random initial
states until convergence to a fixed point or limit cycle (detected by
Floyd's algorithm, maximum 200 steps). The value function is:
%
\begin{equation}
\label{eq:value}
v(S) = \frac{1}{N} \sum_{i=1}^{N}
\overline{x_{\mathrm{out}}^{(i)}(S)},
\end{equation}
%
where $\overline{x_{\mathrm{out}}^{(i)}(S)}$ is the time-averaged
output-node state for initial condition $i$ under coalition $S$. The
global Walsh--Hadamard transform of $v$ gives the \emph{global
interaction spectrum}:
%
\begin{equation}
\label{eq:global-wht}
w_T^{\mathrm{global}} = \frac{1}{2^n} \sum_{S \subseteq \{1,\ldots,n\}}
v(S)\,(-1)^{|S \cap T|}.
\end{equation}
%
The global pairwise magnitude is
$g_{ij} = |w_{\{i,j\}}^{\mathrm{global}}|$, and the
\emph{global energy spectrum} is defined analogously to
Eq.~\ref{eq:local-spectrum}.


\subsection{Composition gap metrics}
\label{sec:metrics}

\textbf{Pairwise Spearman correlation.}
For each network, $\rho = \mathrm{Spearman}(\ell_{ij}, g_{ij})$ over
all $\binom{n}{2}$ pairs, with bootstrap 95\% CIs (10{,}000
resamples).

\textbf{Order-3+ energy gap.}
$\Delta_{3+} = \sum_{m \geq 3} E_m^{\mathrm{global}}
- \sum_{m \geq 3} E_m^{\mathrm{local}}.$
Positive $\Delta_{3+}$ indicates creation; negative indicates
destruction. Networks with $|\Delta_{3+}| < 0.5$~pp are classified as
null.

\textbf{Quality gate.}
The global spectrum is analyzed only if (i)~the value function has
$\geq 10$ unique values, (ii)~absolute order-3+ energy exceeds
$10^{-4}$, and (iii)~no single value accounts for $> 90\%$ of
coalitions. All 27~networks pass.


\subsection{Blind prediction protocol}
\label{sec:blind}

Structural predictions were registered in two batches before coalition
sweeps ran. For Batch~1 (6~networks), predictions were based on
Boolean rule analysis, feedback loop topology, and output-node
structure. A cross-network hypothesis was registered: positive
feedback loop ratio $> 0.5$ implies creation. For Batch~2
(21~networks), predictions used the same structural analysis applied
independently. All predictions were SHA-frozen before any sweep
executed.


%% ====================================================================
\section{Results}
\label{sec:results}
%% ====================================================================

\subsection{Creation is the dominant mode}

Table~\ref{tab:all_results} summarizes all 27 networks. Of 24
non-null networks ($|\Delta_{3+}| \geq 0.5$~pp), 18 show creation
and 6 show destruction. The median $\Delta_{3+}$ is $+8.5$~pp, and
the mean is $+20$~pp. The gap magnitudes span three orders: from
effectively zero (Davidich yeast, $-0.4$~pp; Remy p53, $+0.08$~pp)
to $+83$~pp (cell-cycle transcription network). This creation bias
indicates that nonlinear dynamics more often generate higher-order
epistasis than compress it.

\begin{longtable}{lrrrrl}
\caption{Composition gap results for 27 Boolean GRNs, sorted by
cycling fraction. $\rho$: Spearman correlation between local and
global pairwise magnitudes. $\Delta_{3+}$: global minus local
order-3+ energy (pp). Classification: creation ($> 0.5$~pp),
destruction ($< -0.5$~pp), or null.}
\label{tab:all_results} \\
\toprule
Network & $n$ & $\rho$ & $p$ & $\Delta_{3+}$ & Cycling \\
\midrule
\endfirsthead
\toprule
Network & $n$ & $\rho$ & $p$ & $\Delta_{3+}$ & Cycling \\
\midrule
\endhead
\bottomrule
\endfoot
Arellano root stem & 9 & 0.60 & $< 10^{-4}$ & $+$50.6 & 0.0\% \\
Segment polarity & 11 & 0.21 & 0.126 & $+$19.8 & 0.0\% \\
EMT switch & 12 & 0.38 & 0.002 & $-$3.8 & 0.1\% \\
Lac operon & 13 & 0.08 & 0.492 & $+$47.4 & 0.3\% \\
Cell cycle transcription & 9 & $-$0.38 & 0.021 & $+$83.3 & 0.4\% \\
Morphogenetic checkpoint & 12 & 0.40 & $< 10^{-3}$ & $+$11.9 & 0.4\% \\
Calzone cell fate (full) & 17 & $-$0.02 & 0.855 & $+$58.2 & 0.5\% \\
Mendoza T-helper$^\dagger$ & 13 & 0.21 & 0.059 & $-$0.3 & 1.7\% \\
Saadatpour guard cell & 13 & 0.07 & 0.572 & $+$51.4 & 5.7\% \\
Asymmetric cell division & 9 & $-$0.10 & 0.547 & $+$44.1 & 6.3\% \\
Myeloid progenitors & 11 & 0.51 & $< 10^{-4}$ & $-$2.3 & 6.6\% \\
Tournier apoptosis & 12 & 0.25 & 0.046 & $+$8.5 & 7.4\% \\
Irons cardiac & 15 & 0.63 & $< 10^{-12}$ & $+$10.9 & 9.3\% \\
Za\~nudo T-LGL & 12 & $-$0.19 & 0.120 & $+$33.2 & 10.3\% \\
Hematopoiesis aging & 15 & 0.63 & $< 10^{-12}$ & $+$0.9 & 10.5\% \\
Calzone cell fate (reduced) & 11 & $-$0.07 & 0.586 & $+$37.0 & 11.1\% \\
Li budding yeast & 11 & 0.53 & $< 10^{-4}$ & $+$6.3 & 12.5\% \\
Fumia cell cycle & 14 & 0.27 & 0.011 & $+$52.2 & 12.5\% \\
Remy p53-Mdm2$^\dagger$ & 10 & 0.50 & $< 10^{-3}$ & $+$0.1 & 13.4\% \\
\emph{Drosophila} cell cycle & 14 & 0.12 & 0.257 & $+$6.2 & 14.0\% \\
Lambda phage & 7 & 0.61 & 0.004 & $-$2.5 & 16.0\% \\
\emph{Arabidopsis} cell cycle & 14 & 0.21 & 0.052 & $+$34.3 & 16.4\% \\
Davidich yeast$^\dagger$ & 10 & 0.15 & 0.312 & $-$0.4 & 17.7\% \\
Pair-rule module & 11 & 0.66 & $< 10^{-7}$ & $-$2.7 & 17.8\% \\
Faure cell cycle & 10 & 0.41 & 0.005 & $-$3.0 & 19.6\% \\
Blood stem cell & 11 & 0.59 & $< 10^{-5}$ & $+$4.8 & 24.5\% \\
Fanconi anemia & 15 & 0.31 & 0.002 & $-$5.5 & 47.5\% \\
\midrule
\multicolumn{6}{l}{\footnotesize $^\dagger$Null: $|\Delta_{3+}| < 0.5$~pp.} \\
\multicolumn{6}{l}{\footnotesize 18 creation, 6 destruction, 3 null out of 27 networks.} \\
\end{longtable}


\subsection{Local--global correlation is positive on average but not universal}
\label{sec:rho}

The median Spearman $\rho$ between local and global pairwise
magnitudes is $0.27$ (mean $0.28$). Sixteen of 27 networks have
significant positive correlations ($p < 0.05$). In these networks,
the molecular wiring diagram provides a useful prior for predicting
which gene pairs will show knockout epistasis.

Five networks show negative $\rho$, and one of these
(cell-cycle transcription, $\rho = -0.38$, $p = 0.02$) is
statistically significant. In this network, gene pairs with strong
local Fourier interactions tend to have \emph{weaker} global knockout
epistasis. Negative correlations indicate that dynamical composition
can actively invert the interaction structure encoded in local rules.

The remaining six networks have non-significant positive
correlations, where the molecular wiring diagram carries no detectable
information about knockout epistasis.


\subsection{Limit-cycle prevalence correlates with the composition gap}
\label{sec:cycling}

A negative association between cycling fraction and $\Delta_{3+}$ was
observed post-hoc on the 6 batch-1 networks ($\rho = -0.83$,
$p = 0.04$, but post-hoc so effectively $\sim 0.10$ two-tailed).
This pattern was then tested on the 21 held-out batch-2 networks:
the correlation has the same sign but is not independently
significant ($\rho = -0.37$, $p = 0.10$, $N = 21$). Pooling all 27
networks gives $\rho = -0.51$ ($p = 0.007$), but this combined
statistic includes the discovery sample and should be interpreted
accordingly.

Among non-null networks, the destruction group has a higher median
cycling fraction (16.9\%) than the creation group (8.3\%), but the
distributions overlap substantially and a Mann--Whitney test on the
categorical split is not significant ($U = 30$, $p = 0.12$). The
continuous correlation does not corroborate a clean categorical
separation: counterexamples exist in both directions. EMT switch has
$0.1\%$ cycling but shows destruction ($-3.8$~pp), and blood stem
cell has $24.5\%$ cycling but shows creation ($+4.8$~pp).

A plausible mechanism exists: time-averaging output over limit-cycle
periods smooths coalition-dependent variation in attractor structure,
compressing higher-order Walsh coefficients. Fixed-point attractors
preserve discrete switching behavior, allowing combinatorial
interactions to compound. Cycling fraction does not correlate with
network size ($\rho = 0.09$, $p = 0.87$, $N = 27$), so this
association is not a proxy for the number of genes. The cycling
hypothesis remains suggestive and awaits confirmation on an
independent sample.


\subsection{Structural predictions perform at chance}
\label{sec:blind-results}

The confirmatory test of structural predictions is batch~2, where
predictions were made on 21 networks whose composition gaps had not
been computed. Direction accuracy on batch~2 was 10/21 (48\%),
indistinguishable from chance (Table~\ref{tab:blind}). Predicted
$\rho$ ranges contained the observed value in 6/21 cases (29\%). A
pre-registered cross-network hypothesis (positive feedback loop ratio
$> 0.5$ implies creation) was falsified in both batches.

Batch~1 predictions (4/6 direction accuracy) were developed on the
same 6 networks they are scored against and are therefore circular;
they are reported for completeness but are not evidence of predictive
power.

\begin{table}[t]
\centering
\caption{Blind prediction accuracy. Batch~2 (21 held-out networks) is
the confirmatory test; batch~1 predictions were developed on the same
networks and are circular. Predictions were SHA-frozen before
coalition sweeps ran.}
\label{tab:blind}
\begin{tabular}{lccc}
\toprule
Category & Batch 2 (confirmatory) & Batch 1 (circular) \\
\midrule
Gap direction & 10/21 (48\%) & 4/6 \\
$\rho$ in predicted range & 6/21 (29\%) & 4/6 \\
fb\_ratio hypothesis & Falsified & Falsified \\
\bottomrule
\end{tabular}
\end{table}

This result is itself informative. The composition gap is not
predictable from the structural features most commonly used to
characterize Boolean networks: feedback loop counts, rule complexity,
in-degree distribution, and output-node arity. Predicting whether
dynamics will create or destroy higher-order epistasis requires
information about the full composed dynamical system that cannot be
extracted from local rules or network topology alone.


%% ====================================================================
\section{Discussion}
\label{sec:discussion}
%% ====================================================================

\subsection{Creation dominance}

The most robust finding is asymmetry: dynamics create higher-order
interactions three times more often than they destroy them (18 vs.\ 6
non-null networks). This asymmetry has a natural explanation. Local
Boolean rules have bounded interaction order (at most $r_k$-way,
where $r_k \leq 6$ in these networks), while attractor-level
interactions can involve all $n$ genes through pathway composition.
The ``space'' of possible higher-order interactions grows
combinatorially with $n$ but is sparsely occupied by local rules,
so dynamics have more room to create new interactions than to
suppress existing ones.

\subsection{Implications for knockout screen design}

The molecular wiring diagram is a statistically significant predictor
of pairwise knockout epistasis in 16 of 27 networks (median
$\rho = 0.27$). In these networks, pathway databases provide a
useful prior for prioritizing combinatorial experiments. In five
networks, however, the local--global correlation is negative: the
wiring diagram actively misleads. A knockout screen guided by
molecular interactions in these networks would over-test the
wrong pairs.

The composition gap provides a more specific guide. In creation
networks, dynamics generate interactions absent from any local rule;
a screen guided by molecular interactions alone would miss the
strongest epistatic pairs. In destruction networks, the wiring diagram
over-predicts interactions that dynamics compress.

\subsection{Relationship to prior work}

Gjuvsland et al.\ \citeneeded{Gjuvsland 2007} demonstrated that
statistical epistasis is generic in gene regulatory networks. Our
results extend this in three directions. First, the gap is quantified
order-by-order rather than detected as present/absent. Second, the
gap is signed, with creation as the dominant mode. Third, across 27
networks, the gap is not predictable from standard topological
features of the regulatory graph---a result established by the
failure of pre-registered predictions.

The global epistasis literature \citeneeded{Otwinowski 2018, Reddy
Desai 2021} establishes that nonlinear genotype-phenotype maps
reshape interaction spectra. Our local-vs-global comparison
operationalizes this reshaping as a quantitative spectral
transformation, decomposed by interaction order.


\subsection{Selection bias in published model repositories}

The creation dominance finding (18/6) is the paper's headline result,
and its primary external-validity threat is selection bias.
The networks in this study were drawn from Cell Collective
\citeneeded{Helikar 2012}, PyBoolNet \citeneeded{Klarner 2017}, and
Biodivine---repositories that curate published Boolean models. These
repositories over-represent networks that were published \emph{because}
they exhibit interesting emergent behavior (multistability, complex
attractors, cell-fate decisions). Networks whose dynamics are simple
enough to be uninteresting are less likely to be modeled, published,
and deposited. If complex emergent behavior correlates with creation
of higher-order epistasis---a plausible but untested assumption---then
the 18/6 ratio may overestimate the prevalence of creation in gene
regulatory networks generally. A definitive test would require
constructing networks from interaction databases (such as KEGG or
Reactome) without selecting for dynamical complexity, which is beyond
the scope of this study.

\subsection{Other limitations}

\begin{itemize}
\item Clamp-to-0 baseline only. Clamp-to-1 may yield different gap
  directions, as the asymmetry of Boolean logic (AND vs.\ OR) makes
  the value function baseline-dependent.
\item All genes are players. Excluding constitutive input nodes
  from the player set would change the pairwise space and potentially
  sharpen correlations.
\item Synchronous update only. Asynchronous dynamics would change
  attractor structure and cycling fractions.
\item Sum-aggregation for local pairwise magnitudes
  (Eq.~\ref{eq:local-pair}). Max-aggregation is an untested
  alternative.
\item $n \leq 17$. The exhaustive approach does not scale to
  genome-wide networks. These results characterize the composition
  gap in well-studied regulatory modules, not whole-genome epistasis.
\end{itemize}


%% ====================================================================
\section{Conclusion}
\label{sec:conclusion}
%% ====================================================================

Across 27 Boolean gene regulatory networks, we measured the complete
Walsh interaction spectrum at both the local rule and attractor
levels. Three findings emerge. First, dynamics create higher-order
epistasis three times more often than they destroy it: composition
is a generically constructive process. Second, the molecular wiring
diagram is a useful but unreliable predictor of knockout
epistasis---positive on average, negative in five networks, and at
chance for predicting the gap direction from structural features
alone. Third, limit-cycle prevalence correlates with the composition
gap ($\rho = -0.51$, $p = 0.007$), consistent with oscillatory
averaging as a compression mechanism, though counterexamples prevent
a clean separation.

The complete Walsh spectra, coalition value functions, and
pre-registration artifacts are available at
\textcolor{red}{[URL: repository]}.


%% ====================================================================
%% References
%% ====================================================================

\bibliographystyle{plainnat}
% \bibliography{references}

\end{document}
